% !TEX root =./ECE343_Practice_main.tex
%
% L27 practice set (Current Density, Conductivity, and Resistance): current
% through a surface, Ohm's law in point form in an aluminum wire, the
% resistance of a wire, and the leakage conductance of a coaxial line.

\fullwidth{\textbf{Learning Objectives 6.1--6.3: } I can calculate the total
current through a surface area given a current distribution. I can use Ohm's
Law for electromagnetics to calculate the relationship between a current
distribution and an electric field. I can articulate and identify what makes a
material a good conductor and how conductivity relates to resistance.
\\[4pt]\normalfont\small Show your work. Write a \textbf{Documentation} line at the
end (\textbf{None} if you did not collaborate). Use $e = 1.602\times10^{-19}$ C
and $\epsilon_0 = 8.854\times10^{-12}$ F/m.}

\begin{questions}

%%%%%%%%%%%%%%%%%%%%%%%%%%%%%%%%%%%%%%%%%
\question \textbf{Current through a surface.} A wire of radius $a = 2$ mm lies
along the $z$ axis and carries the current density
$\mathbf{J} = \hat{\mathbf{z}}\ J_0\lp 1 - \rho^2/a^2\rp$, with
$J_0 = 5\times10^{5}\ \text{A/m}^2$.
\begin{parts}

	\part Charge of density $\rho_v$ moves with velocity $\mathbf{u}$. Show
	that the current through a small area $\Delta S$ with unit normal
	$\hat{\mathbf{n}}$ is $\Delta I = \lp\rho_v\mathbf{u}\rp\mycdot\hat{\mathbf{n}}\ \Delta S$.
	\begin{solution}[1.3in]
	In a time $\Delta t$ the charges that cross $\Delta S$ are the ones in the
	tube behind it, of cross section $\Delta S$ and length
	$\mathbf{u}\mycdot\hat{\mathbf{n}}\ \Delta t$ along the normal:
	\eq{
		\Delta Q &= \rho_v\lp\mathbf{u}\mycdot\hat{\mathbf{n}}\ \Delta t\rp\Delta S \\
		\Delta I &= \dfrac{\Delta Q}{\Delta t} = \lp\rho_v\mathbf{u}\rp\mycdot\hat{\mathbf{n}}\ \Delta S
	}
	\end{solution}

	\begin{minipage}[t]{0.58\linewidth}
		\part Find the total current in the wire.
		\begin{solution}[1.5in]
		With $d\mathbf{s} = \hat{\mathbf{z}}\ \rho\ d\rho\ d\phi$:
		\eq{
			I &= 2\pi J_0\int_0^a \lp\rho - \dfrac{\rho^3}{a^2}\rp d\rho \\
			&= 2\pi J_0\lp\dfrac{a^2}{2} - \dfrac{a^2}{4}\rp = \dfrac{\pi a^2 J_0}{2} \\
			&= \dfrac{\pi\times\lp2\times10^{-3}\rp^2\times5\times10^{5}}{2} = 3.14\ \text{A}
		}
		\end{solution}
	\end{minipage}\hfill%
	\begin{minipage}[t]{0.37\linewidth}
		\raggedleft \vspace{12pt}
		$I$ = \ansbox[1.3in]{$3.14\ \text{A}$}
	\end{minipage}

	\begin{minipage}[t]{0.58\linewidth}
		\part Find the current inside $\rho < a/2$, and the fraction of the
		total it carries.
		\begin{solution}[1.5in]
		\eq{
			I_{1/2} &= 2\pi J_0\lp\dfrac{(a/2)^2}{2} - \dfrac{(a/2)^4}{4a^2}\rp
			= \dfrac{7}{32}\pi a^2 J_0 = 1.37\ \text{A} \\
			\dfrac{I_{1/2}}{I} &= \dfrac{7/32}{1/2} = \dfrac{7}{16} = 0.44
		}
		The inner quarter of the area carries $44\%$ of the current, because the
		current density is largest on the axis.
		\end{solution}
	\end{minipage}\hfill%
	\begin{minipage}[t]{0.37\linewidth}
		\raggedleft \vspace{12pt}
		$I_{1/2}$ = \ansbox[1.3in]{$1.37\ \text{A}$}

		\vspace{4pt}
		fraction = \ansbox[1.3in]{$0.44$}
	\end{minipage}

\end{parts}

\newpage
%%%%%%%%%%%%%%%%%%%%%%%%%%%%%%%%%%%%%%%%%
\question \textbf{Ohm's law in point form.} An aluminum wire of cross section
$2.5\ \text{mm}^2$ carries $10$ A uniformly. Aluminum has
$\sigma = 3.7\times10^{7}$ S/m and $n = 1.8\times10^{29}$ free electrons per
cubic meter.
\begin{parts}

	\begin{minipage}[t]{0.58\linewidth}
		\part Find the current density, the field in the wire, and the voltage
		across a $20$ m length.
		\begin{solution}[1.4in]
		\eq{
			J &= \dfrac{I}{A} = \dfrac{10}{2.5\times10^{-6}} = 4.0\times10^{6}\ \text{A/m}^2 \\
			E &= \dfrac{J}{\sigma} = \dfrac{4.0\times10^{6}}{3.7\times10^{7}} = 0.108\ \text{V/m} \\
			V &= E\ell = 0.108\times20 = 2.16\ \text{V}
		}
		\end{solution}
	\end{minipage}\hfill%
	\begin{minipage}[t]{0.37\linewidth}
		\raggedleft \vspace{12pt}
		$J$ = \ansbox[1.3in]{$4.0\times10^{6}$ A/m$^2$}

		\vspace{4pt}
		$E$ = \ansbox[1.3in]{$0.108$ V/m}

		\vspace{4pt}
		$V$ = \ansbox[1.3in]{$2.16$ V}
	\end{minipage}

	\begin{minipage}[t]{0.58\linewidth}
		\part Find the drift speed of the electrons and their mobility.
		\begin{solution}[1.2in]
		\eq{
			u &= \dfrac{J}{ne} = \dfrac{4.0\times10^{6}}{1.8\times10^{29}\times1.602\times10^{-19}} \\
			&= 1.39\times10^{-4}\ \text{m/s} \\
			\mu_e &= \dfrac{\sigma}{ne} = 1.28\times10^{-3}\ \text{m}^2/\text{V}\mycdot\text{s}
		}
		\end{solution}
	\end{minipage}\hfill%
	\begin{minipage}[t]{0.37\linewidth}
		\raggedleft \vspace{12pt}
		$u$ = \ansbox[1.3in]{$0.139$ mm/s}

		\vspace{4pt}
		$\mu_e$ = \ansbox[1.3in]{$1.28\times10^{-3}$}
	\end{minipage}

	\part Aluminum has more free electrons per cubic meter than copper
	($8.5\times10^{28}$), yet a lower conductivity. Explain why, using
	$\sigma = ne\mu_e$.
	\begin{solution}[1in]
	The conductivity is the product of the carrier density and the mobility.
	Aluminum has about twice copper's free-electron density, but its electrons
	collide with the lattice more often, so their mobility is about a third of
	copper's ($1.28\times10^{-3}$ against $4.3\times10^{-3}$), and the product
	is smaller.
	\end{solution}

	\begin{minipage}[t]{0.58\linewidth}
		\part A power line must have the same resistance per meter in aluminum
		as in copper ($\sigma = 5.8\times10^{7}$ S/m). Find the ratio of the
		cross sections, and the ratio of the masses per meter. The densities
		are $2700$ and $8960\ \text{kg/m}^3$.
		\begin{solution}[1.4in]
		$R^{\prime} = 1/(\sigma A)$ is the same, so $\sigma A$ is the same:
		\eq{
			\dfrac{A_\text{Al}}{A_\text{Cu}} &= \dfrac{5.8\times10^{7}}{3.7\times10^{7}} = 1.57 \\
			\dfrac{m_\text{Al}}{m_\text{Cu}} &= 1.57\times\dfrac{2700}{8960} = 0.47
		}
		\end{solution}
	\end{minipage}\hfill%
	\begin{minipage}[t]{0.37\linewidth}
		\raggedleft \vspace{12pt}
		area ratio = \ansbox[1.3in]{$1.57$}

		\vspace{4pt}
		mass ratio = \ansbox[1.3in]{$0.47$}
	\end{minipage}

\end{parts}

\newpage
%%%%%%%%%%%%%%%%%%%%%%%%%%%%%%%%%%%%%%%%%
\question \textbf{Resistance of a wire.} A straight wire of length $\ell$,
cross section $A$, and conductivity $\sigma$ has its end at $z = 0$ held at
voltage $V$ and its end at $z = \ell$ at zero.
\begin{parts}

	\part Assume a uniform field $\mathbf{E} = \hat{\mathbf{z}}\ E$ in the wire.
	Integrate for the voltage and the current, and derive $R = \ell/(\sigma A)$.
	\begin{solution}[1.6in]
	\eq{
		V &= -\int_{\ell}^{0} E\ dz = E\ell \\
		I &= \int_S \sigma E\ \hat{\mathbf{z}}\mycdot\hat{\mathbf{z}}\ ds = \sigma E A \\
		R &= \dfrac{V}{I} = \dfrac{E\ell}{\sigma E A} = \dfrac{\ell}{\sigma A}
	}
	\end{solution}

	\part Explain why the field must have the same magnitude all along the wire
	when the current is steady.
	\begin{solution}[0.9in]
	A steady current has $\nabla\mycdot\mathbf{J} = 0$, so no charge collects
	between two cross sections and the same current crosses each one. Every
	cross section has the same area, so $J$, and $E = J/\sigma$, are the same
	along the wire.
	\end{solution}

	\begin{minipage}[t]{0.58\linewidth}
		\part Find the DC resistance of $100$ m of copper wire of diameter
		$2.05$ mm.
		\begin{solution}[1.1in]
		\eq{
			A &= \pi\lp1.025\times10^{-3}\rp^2 = 3.30\times10^{-6}\ \text{m}^2 \\
			R &= \dfrac{100}{5.8\times10^{7}\times3.30\times10^{-6}} = 0.522\ \ohms
		}
		\end{solution}
	\end{minipage}\hfill%
	\begin{minipage}[t]{0.37\linewidth}
		\raggedleft \vspace{12pt}
		$R$ = \ansbox[1.3in]{$0.522\ \ohms$}
	\end{minipage}

\end{parts}

\newpage
%%%%%%%%%%%%%%%%%%%%%%%%%%%%%%%%%%%%%%%%%
\question \textbf{Leakage through a coaxial line.} A coaxial line has inner
radius $a$, shield radius $b$, and a dielectric of permittivity $\epsilon$ and
small conductivity $\sigma$.
\begin{parts}

	\part Place $\rho_\ell$ per meter on the inner conductor. Use the field from
	Gauss's law to derive the leakage conductance per meter,
	$G^{\prime} = 2\pi\sigma/\ln(b/a)$.
	\begin{solution}[2in]
	\eq{
		\mathbf{E} &= \hat{\boldsymbol{\rho}}\dfrac{\rho_\ell}{2\pi\epsilon\rho}, \qquad
		V = -\int_b^a \dfrac{\rho_\ell}{2\pi\epsilon\rho}\ d\rho = \dfrac{\rho_\ell}{2\pi\epsilon}\ln\dfrac{b}{a} \\
		I &= \int_0^{\ell}\int_0^{2\pi}\sigma\dfrac{\rho_\ell}{2\pi\epsilon\rho}\ \rho\ d\phi\ dz = \dfrac{\sigma\rho_\ell\ell}{\epsilon} \\
		G &= \dfrac{I}{V} = \dfrac{2\pi\sigma\ell}{\ln(b/a)}, \qquad G^{\prime} = \dfrac{G}{\ell} = \dfrac{2\pi\sigma}{\ln(b/a)}
	}
	\end{solution}

	\begin{minipage}[t]{0.58\linewidth}
		\part Moisture raises the conductivity of the dielectric in a $50$ m run
		of RG-58 ($a = 0.405$ mm, $b = 1.475$ mm, $\epsilon_r = 2.25$) to
		$10^{-12}$ S/m. Find $G^{\prime}$ and the DC resistance between the
		conductors.
		\begin{solution}[1.3in]
		\eq{
			G^{\prime} &= \dfrac{2\pi\times10^{-12}}{\ln(1.475/0.405)} = \dfrac{6.283\times10^{-12}}{1.293} \\
			&= 4.86\times10^{-12}\ \text{S/m} \\
			R &= \dfrac{1}{G^{\prime}\ell} = \dfrac{1}{4.86\times10^{-12}\times50} = 4.1\ \text{G}\ohms
		}
		\end{solution}
	\end{minipage}\hfill%
	\begin{minipage}[t]{0.37\linewidth}
		\raggedleft \vspace{12pt}
		$G^{\prime}$ = \ansbox[1.3in]{$4.86$ pS/m}

		\vspace{4pt}
		$R$ = \ansbox[1.3in]{$4.1\ \text{G}\ohms$}
	\end{minipage}

	\part Lesson 25 gives $C^{\prime} = 96.8$ pF/m for this cable. Show that
	$G^{\prime}/C^{\prime} = \sigma/\epsilon$ with your numbers.
	\begin{solution}[1in]
	\eq{
		\dfrac{G^{\prime}}{C^{\prime}} &= \dfrac{4.86\times10^{-12}}{96.8\times10^{-12}} = 0.0502\ \text{s}^{-1} \\
		\dfrac{\sigma}{\epsilon} &= \dfrac{10^{-12}}{2.25\times8.854\times10^{-12}} = 0.0502\ \text{s}^{-1}
	}
	The logarithm that carries the geometry cancels in the ratio.
	\end{solution}

\end{parts}

\vspace{1.5em}
\noindent\textbf{Documentation:} \rule[-2pt]{0.6\linewidth}{0.4pt}

\end{questions}
